\chapter{The domain}
\label{ch:domain}

\section{A network as a heterogeneous directed graph}
\begin{definition}[Network graph]
A network is a tuple $G=(\V,\E,\tau,\rho)$ where $\V$ is a finite set of $N$ nodes,
$\E\subseteq\V\times\V$ a set of directed edges, $\tau:\V\to\Tset$ a node-type map and
$\rho:\E\to\Rel$ an edge-relation map.
\begin{equation}
G=(\V,\E,\tau,\rho),\qquad
\Tset=\{\textsf{cell},\textsf{gnb},\textsf{router},\textsf{link},\textsf{upf},\textsf{amf},\textsf{smf},\textsf{service}\}.
\label{eq:graph}
\end{equation}
\end{definition}

Three choices in \eqref{eq:graph} matter and are easy to get wrong.

\paragraph{Edges point in the direction a fault propagates.} An edge $v\to u$ means ``a
degradation at $v$ is felt at $u$''. A fibre link feeds a router; a core router feeds an aggregation
router; an aggregation router backhauls a gNB; a gNB serves cells; the AMF controls gNBs; the SMF
controls UPFs; a core router feeds a UPF; a UPF terminates services; cells carry services. This
is the \emph{causal} orientation, not the traffic orientation, and it is what makes
Chapter~\ref{ch:propagation} a diffusion and Chapter~\ref{ch:inverse} an inverse problem.

\paragraph{Links are nodes.} A transport link carries its own KPIs (latency, loss, utilisation)
and can be the origin of a fault. Modelling it as an edge attribute would make it impossible to
name as a root cause. So links are first-class nodes of type \textsf{link} with a single outgoing
relation to the router they feed.

\paragraph{Relations are typed.} $\Rel$ has nine members. ``Link feeds router'' and ``AMF controls
gNB'' must not share a weight matrix: a signalling storm and a fibre cut leave different
signatures on their neighbours.

\begin{figure}[h]
\centering
\begin{tikzpicture}[
  >=Latex, node distance=9mm,
  n/.style={draw, rounded corners=2pt, minimum width=14mm, minimum height=6.5mm, font=\small, thick},
  ran/.style={n, draw=ran, fill=ran!8}, tra/.style={n, draw=tra, fill=tra!8},
  core/.style={n, draw=core, fill=core!10}, svc/.style={n, draw=svc, fill=svc!8},
  e/.style={->, thick, draw=ink!70}
]
\node[core] (amf) {amf};
\node[core, right=of amf] (smf) {smf};
\node[core, below=of smf] (upf) {upf};
\node[tra, left=22mm of upf] (rc) {router (core)};
\node[tra, below=of rc] (ra) {router (agg)};
\node[tra, left=of ra] (lk) {link};
\node[ran, below=of ra] (gnb) {gnb};
\node[ran, below=of gnb] (cell) {cell};
\node[svc, right=22mm of cell] (s) {service};
\draw[e] (lk) -- (ra) node[midway, above, font=\scriptsize] {feeds};
\draw[e] (rc) -- (ra) node[midway, right, font=\scriptsize] {feeds};
\draw[e] (ra) -- (gnb) node[midway, right, font=\scriptsize] {backhauls};
\draw[e] (gnb) -- (cell) node[midway, right, font=\scriptsize] {serves};
\draw[e] (amf) to[out=200, in=160] node[pos=0.5, left, font=\scriptsize] {controls} (gnb);
\draw[e] (smf) -- (upf) node[midway, right, font=\scriptsize] {controls};
\draw[e] (rc) -- (upf) node[midway, above, font=\scriptsize] {feeds};
\draw[e] (upf) -- (s) node[midway, right, font=\scriptsize] {terminates};
\draw[e] (cell) -- (s) node[midway, above, font=\scriptsize] {carries};
\begin{scope}[on background layer]
\node[draw=ran!60, dashed, rounded corners, fit=(gnb)(cell), inner sep=4pt, label={[ran, font=\scriptsize]left:RAN}] {};
\node[draw=tra!60, dashed, rounded corners, fit=(rc)(ra)(lk), inner sep=4pt, label={[tra, font=\scriptsize]left:Transport}] {};
\node[draw=core!70, dashed, rounded corners, fit=(amf)(smf)(upf), inner sep=4pt, label={[core, font=\scriptsize]above:Core}] {};
\end{scope}
\end{tikzpicture}
\caption{The node types and relations of \eqref{eq:graph}. Arrows point in the direction a fault
propagates. A service is the trigger surface: a service-layer symptom is the only input the
healing loop needs.}
\label{fig:domain}
\end{figure}

\section{Adjacency and signals}
Let $A\in\{0,1\}^{N\times N}$ be the cause$\to$effect adjacency and $A_r$ its restriction to
relation $r\in\Rel$:
\begin{equation}
A_{vu}=\ind{(v,u)\in\E},\qquad A=\sum_{r\in\Rel}A_r,\qquad
\N_r(v)=\{u:(u,v)\in\E,\ \rho(u,v)=r\}.
\label{eq:adjacency}
\end{equation}
Each node carries a KPI vector $x_v(t)\in\R^K$ at discrete time $t$, standardised so that nominal
behaviour is near zero. We use $K=4$ (throughput, latency, loss, utilisation) for every type; the
tensors stay rectangular and the \emph{direction} in which degradation shows up is what differs by
type (Chapter~\ref{ch:propagation}). The full state is a signal on the graph,
\begin{equation}
X(t)\in\R^{N\times K},\qquad X_{v\cdot}(t)=x_v(t).
\label{eq:signal}
\end{equation}

\section{Business value as a node attribute}
Every cell $c$ carries a revenue density $\mathrm{rev}_c\ge 0$. It is heavy-tailed on purpose
(log-normal in the generator): a few cells carry most of the revenue, which is the regime where
``same fault, different response'' is not a slogan but a measurable reordering
(Chapter~\ref{ch:intent}). Revenue is an attribute of the twin, not of the model: the model never
sees it, and that separation is the whole design.
